\section{Introduction}
\label{sec:intro}

Continuous water-quality monitoring underpins the management of rivers that supply drinking water and receive pollution. In-situ sensors measure temperature, specific conductance, dissolved oxygen, and pH every few minutes, and the resulting records reveal diel cycles, storm responses, and short pulses that grab samples miss \citep{rode2016}. Wireless sensor networks have made such monitoring affordable \citep{adumanu2017}, and deployments in Ghana and Kosovo have shown that it works in practice, but they relied on cellular backhaul \citep{adumanu2020,ahmedi2022}. The rivers under the greatest pressure are often the furthest from that infrastructure; in Ghana, for example, small-scale gold mining has degraded rivers in forested catchments where terrestrial coverage is sparse \citep{duncan2020}. This motivating case is not evaluated here.

Satellite Internet of Things (IoT) services reach such sites with small, low-power terminals \citep{centenaro2021}, but their capacity is sold in messages, not bandwidth. A typical module sends one-way messages with a payload of 23 bytes and receives no acknowledgement \citep{kineis2023kim1}. A four-parameter sensor sampling every 10 minutes produces 576 readings a day; even packed at 13 bits each, without timestamps, into the 20 bytes that remain after a three-byte message header, they would fill at least 47 messages. Few such messages are affordable, some are lost, and the node is never told which. A monitoring agency therefore needs to know what the server can still conclude from what arrived. Many uses tolerate a delay of a day: compliance assessment, trend analysis, and model calibration work from daily or longer summaries, whereas alarms that require action within minutes are outside the scope of this study.

Agencies do not act on records; they act on decisions such as whether dissolved oxygen fell below 5 mg/L or whether a conductance pulse occurred. Delivering a whole record within an accuracy tolerance \citep{wagner2006} serves such decisions only partly. A record within tolerance cannot settle a reading that lies closer to a threshold than the tolerance allows, and when the quota is binding, bytes spent on the whole record are bytes not spent near the threshold. This study therefore took the decision as the object: the server states, for each day and criterion, one of three outcomes (exceeded, not exceeded, or undetermined), and the encoder spends its bytes so that the outcomes that matter are decided, with soundness that holds whatever the loss pattern, provided the messages that arrive are uncorrupted.

The study asked five research questions:
\begin{enumerate}
  \item RQ1: Under which assumptions is a three-valued threshold decision from received bytes sound, and what can no decoder decide?
  \item RQ2: Do compact extremum witnesses, and targeted protection of the witnesses of hard cases, raise the share of decided hard cases over a record alone at a fixed byte quota?
  \item RQ3: How does the witness design compare with decision flags (the answers to configured thresholds sent as bits), send-on-delta sampling, and a tightened tolerance, including on thresholds that were not known when the node was configured?
  \item RQ4: What does it cost, in record recovery and in message capacity?
  \item RQ5: How much does a message that is corrupted but passes its check endanger the soundness of the decisions?
\end{enumerate}

Five terms recur throughout:
\begin{itemize}
  \item Certificate: the server's three-valued answer to a threshold query for one day (exceeded, not exceeded, or undetermined), computed only from the bytes that arrived; a certain answer (exceeded or not exceeded) is correct under the conditions stated in Section~\ref{sec:model}.
  \item Witness: a compact value sent beside the record that settles a decision the record alone may not, such as the quantised daily minimum of dissolved oxygen.
  \item Hard case: a day on which the record alone may not settle a query, because the relevant value lies near the threshold or the threshold is crossed; the bands that define it are given in Section~\ref{sec:model}.
  \item Configured and unconfigured thresholds: configured thresholds are known to the node when it is set up; unconfigured thresholds are thresholds the node was not set up for, such as a criterion revised after deployment. In this study the unconfigured thresholds were chosen by the analysts and are close to the configured ones.
  \item Whole-record recovery: the share of days on which every planned value of the error-bounded record arrived; on such days every reading is certified within tolerance.
\end{itemize}

The broader principle is that a severely bandwidth-limited monitoring system should keep the information that the decisions of a future receiver need, rather than only the information that reconstructs the record; the study tests one implementation of that principle for four water-quality parameters.

The building blocks are established, and the study does not claim otherwise: minimum-size piecewise-linear approximation within a maximum error \citep{imai1988,orourke1981}; budgeted compression of sensor history \citep{deligiannakis2007}; priority encoding transmission, which protects parts of a message unequally \citep{albanese1996,boucheron2000pet}; deterministic interval answers from compressed or uncertain data \citep{cheng2003probabilistic,brito2017platodb,lin2019plato}; possible-minimum and possible-maximum filtering of predicates \citep{taurone2024treat}; and send-on-delta sampling \citep{diazcacho2017sod}. The contributions of this study are the following:
\begin{itemize}
  \item A certificate language for a one-way lossy link: every piece of received information becomes an interval on a reading; threshold, extremum, conductance-change, and rolling-mean queries receive three-valued answers; the answers are sound when the intervals are valid (proved and checked exhaustively on small cases); and two limits are stated, namely that the rules are exact only without witnesses and that the tolerance makes near-threshold readings undecidable from a record alone.
  \item A message frame with one CRC-24 per transmitted message, a witness section of 13-bit codes, and a hardened receiver (a sanity rule, corroboration, and rejection of any block whose decoding fails), together with a byte-level codec, fuzz tests, and a mutation check of the tests.
  \item A comparison of designs on identical loss draws at byte level: record only, targeted witnesses (the 3b rule), and a fair flags comparator that sends the answers to the configured thresholds and also uses the record; in development, an exact allocator, send-on-delta, a tightened tolerance, and two hybrids, including negative results (the exact allocator did not beat the simple rule, and flags win on configured thresholds).
  \item Two pre-registered confirmations on sealed monitoring sites in years that development had not used. The first (version 2 of the method, 31 sites, 2024) failed its integrity hypothesis, and an internal review then found that its flags comparator discarded the record; the method was revised (version 3), and the earlier cohort became development data. The second (version 3, 35 sites, 2022) tested three hypotheses that could fail; all held. A modelling error found before the second freeze, and a design adopted although it missed its own adoption limit, are reported as they occurred.
\end{itemize}

Section~\ref{sec:related} reviews related work. Section~\ref{sec:model} describes the method. Section~\ref{sec:study} describes the data and the pre-registered design. Section~\ref{sec:results} reports the development and confirmation results, Section~\ref{sec:discussion} discusses them and their limits, and Section~\ref{sec:conclusion} concludes.
